% Source: https://artofproblemsolving.com/community/c3h2232214_proof_of_quadratic_formula_2
% Stable identity: aops-topic-2232214-q01
\textbf{(a)} Multiply $ax^2+bx+c=0$ by $4a\ne0$ (this gives an equivalent equation) and complete the square:
\[4a^2x^2+4abx+4ac=0\iff(2ax+b)^2=b^2-4ac .\]
If $b^2-4ac\ge0$, this holds exactly when $2ax+b=\pm\sqrt{b^2-4ac}$, i.e. when
\[x=\dfrac{-b\pm\sqrt{b^2-4ac}}{2a}.\]

\textbf{(b)} With $b=2b'$ we have $b^2-4ac=4\bigl((b')^2-ac\bigr)\ge0$, so by (a)
\[x=\dfrac{-2b'\pm\sqrt{4\bigl((b')^2-ac\bigr)}}{2a}=\dfrac{-2b'\pm2\sqrt{(b')^2-ac}}{2a}=\dfrac{-b'\pm\sqrt{(b')^2-ac}}{a}.\]
(Directly: multiplying $ax^2+2b'x+c=0$ by $a$ gives $(ax+b')^2=(b')^2-ac$.) $\blacksquare$

For example, $x^2+6x-7=0$ has $b'=3$ and $x=-3\pm\sqrt{9+7}=-3\pm4$, i.e. $x=1$ or $x=-7$.
